\documentclass[11pt]{article}
\usepackage[a4paper,margin=1.8cm]{geometry}
\usepackage{amsmath, amssymb, amsthm}  % Essential math packages
\usepackage{graphicx}                   % For figures
\usepackage{hyperref}                   % Clickable links
\usepackage{parskip} % cleaner spacing
\usepackage{enumitem}
\usepackage{amssymb}

\title{Sections and Chapters}

\newtheorem{definition}{Definisjon}

\title{Oblig 1 \\ MAT1100 - Lineær Algebra}
\author{Oliver Ekeberg}
\date{\today}

\begin{document}
\maketitle


\tableofcontents


\section{2.5}

\subsection{oppgave 2.5.2}

Om vi har alle n-dim vektor rom over $ \mathbb{K} $, så finnes det en $ T \in \mathcal{L}(U, V) $  


\subsection{oppgave 2.5.3}

Har at $ \mathcal{P}_2 $ har basis $ \mathcal{B} = (p_0, p_1,p_2)$ den kanoniske basisen og $ \mathcal{C}=(q_0,q_1,q_2) $ newton basisen over punktene $ t_0=0,t_1=1,t_2=2 $

hvor $ p(t)=2-t+3t^2 $ 

Da er 

\begin{align*}
    \left[ p \right]_{\mathcal{B}} &=2p_0 + (-1)p_1+3p_2 \\
    &= (2,-1,3)
\end{align*}




\end{document}
